\chapter{Quick Reference and Tables}
łabel{ch:gm−reference}
∈dex{quick reference}

This chapter collects quick−reference tables, optional rules, and GM tools. Use them as needed at the table. For additional player−facing reference, see the \textit{Player's Guide} appendices.

\clearpage
\section*{Appendix A: Core Reference Tables}
\addcontentsline{toc}{section}{Appendix A: Core Reference Tables}
łabel{app:core−tables}

\begin{longtable}{lll}
∩tion{Outcome Matrix} łabel{tab:outcome−matrix} \\
⊤rule
\textbf{Case} & \textbf{Name} & \textbf{Guidance} \\
∣rule
\endfirsthead
μlticolumn{3}{c}%
{{\tablename\ θble{} −− continued}} \\
⊤rule
\textbf{Case} & \textbf{Name} & \textbf{Guidance} \\
∣rule
\endhead
⊥tomrule
\endfoot
$S ≥ DV$, $C = 0$ & Clean Success & Deliver the intent crisply. \\
$S ≥ DV$, $C > 0$ & Success & Cost & Grant intent; spend or bank SB for complications. \\
$0 < S < DV$ & Partial & Progress; award 1 Boon. \\
$S = 0$ & Miss & No progress; award 2 Boons; spend or bank SB. \\
\end{longtable}
∈dex{outcome matrix}∈dex{clean success}∈dex{success with sb}∈dex{partial success}∈dex{miss}

\begin{tcolorbox}[title=Story Beat (SB) Spend Menu, colback=white!97!gray, colframe=black!80!gray]
∈dex{story beats!spending}
\begin{itemize}
    ℹtem \textbf{1 SB}: Minor pressure −− noise, trace, \(+1\) segment on a relevant timer.
    ℹtem \textbf{2 SB}: Moderate setback −− alarm raised, lose cover or advantage, lesser foe appears.
    ℹtem \textbf{3 SB}: Serious trouble −− reinforcements arrive, key gear breaks temporarily, tick a campaign front.
    ℹtem \textbf{4+ SB}: Major turn −− trap springs, authority arrives, scene shifts dramatically.
\end{itemize}
\end{tcolorbox}

\begin{tcolorbox}[title=Position Descriptions, colback=white!97!gray, colframe=black!80!gray]
∈dex{position!descriptions}
\begin{itemize}
    ℹtem \textbf{Dominant}: Act on your terms. Consequences are manageable.
    ℹtem \textbf{Controlled}: Normal risk and reward. Failure has a real cost.
    ℹtem \textbf{Desperate}: Odds stacked against you. Failure could be catastrophic.
\end{itemize}
\end{tcolorbox}

\begin{longtable}{cll}
∩tion{Difficulty Ladder} łabel{tab:dv−ladder} \\
⊤rule
\textbf{DV} & \textbf{Name} & \textbf{When to Use} \\
∣rule
\endfirsthead
μlticolumn{3}{c}%
{{\tablename\ θble{} −− continued}} \\
⊤rule
\textbf{DV} & \textbf{Name} & \textbf{When to Use} \\
∣rule
\endhead
⊥tomrule
\endfoot
2 & Routine & Clear intent, modest stakes, controlled environment. \\
3 & Pressured & Time pressure, mild opposition, partial information. \\
4 & Hard & Hostile conditions, active opposition, precise timing. \\
5+ & Extreme & Multiple constraints, high stakes, dramatic failure. \\
\end{longtable}
∈dex{difficulty value!ladder}

\begin{tcolorbox}[title=Boon Economy Quick Guide, colback=white!97!gray, colframe=black!80!gray]
∈dex{boons!economy}
\begin{itemize}
  ℹtem \textbf{Holding cap:} 5 Boons maximum.
  ℹtem \textbf{Carryover:} At the end of each scene, reduce held Boons to a maximum of 2.
  ℹtem \textbf{Conversion:} 2 Boons → 1 XP, once per session during downtime (max 2 XP per session via conversion).
  ℹtem \textbf{Spending:} Re−roll a die, activate an asset, improve Position, or power certain Rites.
\end{itemize}
\end{tcolorbox}

\begin{tcolorbox}[title=Mechanical Constraints, colback=white!97!gray, colframe=black!80!gray]
∈dex{assist limit}∈dex{follower actions}
\begin{itemize}
  ℹtem \textbf{Assist Max:} \(+3\) dice total per roll (Exceptional Coordination allows \(+4\) from one follower).
  ℹtem \textbf{Boon Max:} 5 total; 2→1 XP conversion once per session (max 2 XP via conversion).
  ℹtem \textbf{Follower Independent Action:} Once per scene (party−wide).
  ℹtem \textbf{Over−Stack:} Two or more structural advantages on the same action → start one obstacle at \(+1\) DV or GM banks \(+1\) SB.
\end{itemize}
\end{tcolorbox}

\begin{longtable}{llll}
∩tion{Harm Levels and Effects} łabel{tab:harm} \\
⊤rule
\textbf{Level} & \textbf{SB Generation} & \textbf{Penalty} & \textbf{Recovery} \\
∣rule
\endfirsthead
μlticolumn{4}{c}%
{{\tablename\ θble{} −− continued}} \\
⊤rule
\textbf{Level} & \textbf{SB Generation} & \textbf{Penalty} & \textbf{Recovery} \\
∣rule
\endhead
⊥tomrule
\endfoot
Minor & 1 SB (next 2 rolls) & −−1 die (related actions) & Rest or basic care \\
Moderate & 1 SB (next roll) & −−1 die (most actions) & Medical treatment \\
Severe & 2 SB (next roll) & −−2 dice (most actions) & Extended care \\
Critical & 3 SB (next roll) & Incapacitated & Major intervention \\
\end{longtable}
∈dex{harm!levels}

\begin{tcolorbox}[title=Range Bands, colback=white!97!gray, colframe=black!80!gray]
∈dex{range bands}
\begin{description}
ℹtem[Close] Arm's length, grappling distance.
ℹtem[Near] Same room or immediate area.
ℹtem[Far] Visible but not immediately reachable.
ℹtem[Absent] Off−screen; requires a scene change.
\end{description}
\end{tcolorbox}

\begin{tcolorbox}[title=Movement Actions, colback=white!97!gray, colframe=black!80!gray]
∈dex{movement actions}
\begin{itemize}
  ℹtem \textbf{Move}: Shift one range band (Close\(\(↔\)\)Near or Near\(\(↔\)\)Far).
  ℹtem \textbf{Dash}: Shift two bands as your full action.
  ℹtem \textbf{Disengage}: Test to leave an enemy's Close range without provoking an opportunity attack.
  ℹtem \textbf{Charge / Skirmish}: Full−turn actions that move two bands and attack, costing 1 Fatigue.
\end{itemize}
\end{tcolorbox}

\begin{tcolorbox}[title=Supply Timer States, colback=white!97!gray, colframe=black!80!gray]
∈dex{supply timer}
\begin{description}
ℹtem[Full (0 segments)] Well−provisioned, no penalty.
ℹtem[Low (2 segments)] Minor friction −− resource checks at −−1 die.
ℹtem[Dangerously Low (3 segments)] Each PC gains Fatigue 1.
ℹtem[Empty (4 segments)] Severe penalties −− starvation or thirst risk.
\end{description}
\end{tcolorbox}

\begin{tcolorbox}[title=Travel Timer Sizes, colback=white!97!gray, colframe=black!80!gray]
∈dex{travel timers}
\begin{description}
ℹtem[4 segments] Short, straightforward leg.
ℹtem[6 segments] Standard leg.
ℹtem[8 segments] Extended or complex journey.
ℹtem[10 segments] Epic or highly dangerous passage.
\end{description}
\end{tcolorbox}

\begin{longtable}{lll}
∩tion{XP Costs} łabel{tab:xp−costs} \\
⊤rule
\textbf{Improvement} & \textbf{Cost} & \textbf{Downtime} \\
∣rule
\endfirsthead
μlticolumn{3}{c}%
{{\tablename\ θble{} −− continued}} \\
⊤rule
\textbf{Improvement} & \textbf{Cost} & \textbf{Downtime} \\
∣rule
\endhead
⊥tomrule
\endfoot
Attribute increase & new rating × 3 XP & new rating days \\
Skill increase & new level × 2 XP & new level days \\
On−Screen Follower & Cap² XP & 1−−3 days \\
Minor Asset & 4 XP & 1 day \\
Standard Asset & 8 XP & 1 week \\
Major Asset & 12 XP & 1 month \\
\end{longtable}
∈dex{XP!costs}

\begin{tcolorbox}[title=Upkeep (Condensed), colback=white!97!gray, colframe=black!80!gray]
∈dex{upkeep}
\textbf{Frequency:} Once per downtime period.
\begin{itemize}
  ℹtem \textbf{Efficient (higher XP, less time):} Upkeep XP = $\max(1, \frac{\text{Acquisition XP}}{3})$.
  ℹtem \textbf{Intensive (1 XP, dedicated action):} One full downtime action.
  ℹtem \textbf{Failure to pay:} Follower becomes Wary (or Seized if already Wary); Asset becomes Neglected (or Compromised if already Neglected).
\end{itemize}
\end{tcolorbox}

\clearpage
\section*{Appendix B: Optional Rules and Tactical Tools}
\addcontentsline{toc}{section}{Appendix B: Optional Rules and Tactical Tools}
łabel{app:optional−rules}

\begin{tcolorbox}[title=Grid−Based Combat (Optional), colback=white!97!gray, colframe=black!80!gray]
∈dex{combat!grid}
Use a square or hex grid with Zones of Control (ZoC).

\textbf{ZoC rules:}
\begin{itemize}
  ℹtem Entering an enemy's Zone of Control ends your movement −− you are engaged.
  ℹtem You cannot move through an enemy's ZoC.
  ℹtem Leaving a ZoC requires the Disengage action (DV 4−−6) or spending 1 Boon.
  ℹtem Flank: \(+1\) die to attack. Rear: \(+1\) die and \(+1\) Effect.
\end{itemize}
\textbf{Push−Through:} Roll $\text{Body}+\text{Athletics}$ vs. DV $= 2 + \text{enemy Tier}$. Choose a cost on success: Position worsens by one step, −−1 die on the action, or the enemy gets a free opportunity attack.
\end{tcolorbox}

\begin{tcolorbox}[title=Detailed Warfare (Mass Combat), colback=white!97!gray, colframe=black!80!gray]
∈dex{combat!mass}
Treat an army as a high−Cap Follower (see \ref{sec:followers−assets−gm}) with its own timers:
\begin{itemize}
  ℹtem \textbf{Supply Timer } (6−−8 segments): When empty, the army suffers severe penalties.
  ℹtem \textbf{Morale Timer } (6−−8 segments): When full, the army routs.
  ℹtem \textbf{Command rolls} (Presence + Command or Wits + Tactics) advance objective timers (e.g., \emph{Break Their Line} [6]).
  ℹtem \textbf{Battle's Edge:} Story Beats generated by command rolls can be spent on fog−of−war complications −− false orders, broken supply wagons, or a messenger's death.
\end{itemize}
\end{tcolorbox}

\begin{tcolorbox}[title=Hex−Based Exploration, colback=white!97!gray, colframe=black!80!gray]
∈dex{exploration!hex}
Map a region in 6−mile hexes. Entering a hex requires a $\text{Wits}+\text{Survival}$ roll (DV set by terrain). Outcomes determine discovery (find a point of interest), danger (ambush, hazard), or becoming lost (add a timer segment to find the way out).
\end{tcolorbox}

\begin{tcolorbox}[title=Miniatures Mode −−− Speed Defaults, colback=white!98!gray, colframe=black!80!gray]
∈dex{miniatures}
\textbf{Base DV:} $\mathrm{DV}=\mathrm{Tier}+2+\text{Keywords}$ (elevation \(+1\), \texttt{[WARD]} \(+1\), Disengage base DV 4).\\
\textbf{Critical Hit (10):} Bump Position one step; if already Dominant, Push or Pull the target 1 hex \emph{or} gain \(+1\) automatic success.\\
\textbf{Zone of Control (ZoC):} Entering or leaving a hex adjacent to an enemy provokes one Reaction (Free Strike or Shove 1 hex).\\
\textbf{Heat:} On a critical hit, the GM may spend 1 Heat (a resource tracked per scene) to degrade Position or trigger a terrain hazard.
\end{tcolorbox}

\begin{tcolorbox}[title=Hex Keywords, colback=white!98!gray, colframe=black!80!gray]
∈dex{hex keywords}
\textbf{Difficult:} Costs 2 movement points per hex. 
\textbf{Elevation:} Attacks from lower ground suffer \(+1\) DV. 
\textbf{ZoC:} Provokes a Reaction on crossing.\\
\textbf{Altar \texttt{[WARD]}:} −−1 die to attacks (or attacker may spend 2 Fatigue to ignore).\\
\textbf{Incorporeal:} Ignores Difficult terrain; can pass through occupied hexes but cannot end movement in one.
\end{tcolorbox}

\subsection*{Sample NPCs}
łabel{sec:sample−npcs}

\begin{tcolorbox}[title=Encounters (Tier I−−II), colback=white!97!gray]
\begin{itemize}
  ℹtem \textbf{Bandit Skirmisher:} Body 2, Wits 2. Melee 2, Stealth 1.
  ℹtem \textbf{Ykrul Rider:} Body 4, Wits 3. Riding 3, Melee 3.
  ℹtem \textbf{Street Bravo:} Presence 3, Body 2. Dueling 3.
\end{itemize}
\end{tcolorbox}

\begin{tcolorbox}[title=Foils & Rivals (Tier II−−III), colback=white!97!gray]
\begin{itemize}
  ℹtem \textbf{Ambitious Scribe:} Wits 3, Presence 3. Intrigue 3, Lore 2.
  ℹtem \textbf{Mercenary Captain:} Body 4, Spirit 3. Command 3, Melee 4.
  ℹtem \textbf{Flame Preacher:} Presence 4, Spirit 3. Oratory 4, Faith 3.
\end{itemize}
\end{tcolorbox}

\begin{tcolorbox}[title=Prestige NPCs (Tier IV−−V), colback=white!97!gray]
\begin{itemize}
  ℹtem \textbf{High Elf Loremaster:} Wits 5, Spirit 4. Lore 5, Arcana 4.
  ℹtem \textbf{Dwarven Forge−Patriarch:} Body 5, Spirit 4. Craft 5, Command 4.
  ℹtem \textbf{Ykrul Warglord:} Body 5, Presence 4. Command 4, Melee 5.
\end{itemize}
\end{tcolorbox}

\begin{tcolorbox}[title=Common Rolls (GM Screen Reference), colback=white!97!gray, colframe=black!80!gray]
∈dex{skill rolls!common examples}
\begin{multicols}{2}
\textbf{Athletics:} climb (\emph{Body}+Athletics), time a leap (\emph{Wits}+Athletics)\\
\textbf{Stealth:} shadow patrol (\emph{Wits}+Stealth), hold still (\emph{Spirit}+Stealth)\\
\textbf{Endurance:} resist cold (\emph{Spirit}+Endurance), push through pain (\emph{Spirit}+Endurance)\\
\textbf{Craft:} blueprint (\emph{Wits}+Craft), brace a door (\emph{Body}+Craft)\\
\textbf{Melee:} break guard (\emph{Body}+Melee), bind a blade (\emph{Wits}+Melee)\\
\textbf{Ranged:} leading shot (\emph{Wits}+Ranged), loose in a squall (\emph{Spirit}+Ranged)\\
\textbf{Brawl:} grapple (\emph{Body}+Brawl), feint (\emph{Wits}+Brawl)\\
\textbf{Tactics:} ambush lane (\emph{Wits}+Tactics), coordinate a retreat (\emph{Presence}+Tactics)\\
\textbf{Diplomacy:} audience (\emph{Presence}+Diplomacy), draft terms (\emph{Wits}+Diplomacy)\\
\textbf{Sway:} haggle (\emph{Presence}+Sway), sell a risky plan (\emph{Presence}+Sway)\\
\textbf{Deception:} lie under scrutiny (\emph{Presence}+Deception), misdirect guard (\emph{Wits}+Deception)\\
\textbf{Performance:} captivate a crowd (\emph{Presence}+Performance), mimic voice (\emph{Wits}+Performance)\\
\textbf{Subterfuge:} talk past a checkpoint (\emph{Presence}+Subterfuge), case a routine (\emph{Wits}+Subterfuge)\\
\textbf{Streetwise:} find a fence (\emph{Presence}+Streetwise), sift rumors (\emph{Wits}+Streetwise)\\
\textbf{Arcana:} read a ward anchor (\emph{Wits}+Arcana), hold a rite steady (\emph{Spirit}+Arcana)\\
\textbf{Mechanics:} diagnose a lock (\emph{Wits}+Mechanics), field−rig a pump (\emph{Body}+Mechanics)\\
\textbf{Investigation:} reconstruct a scene (\emph{Wits}+Investigation), interview a witness (\emph{Presence}+Investigation)\\
\textbf{Lore:} cite local custom (\emph{Presence}+Lore), perform a rite (\emph{Spirit}+Lore)\\
\textbf{Nature:} read weather (\emph{Wits}+Nature), calm a mount (\emph{Presence}+Nature)\\
\textbf{Medicine:} stabilize a wound (\emph{Wits}+Medicine), cut away rot (\emph{Body}+Medicine)\\
\textbf{Command:} rally allies (\emph{Presence}+Command), hold the line (\emph{Spirit}+Command)
\end{multicols}
\vspace{0.5em}
\textbf{Fast Boundaries:}
\begin{itemize}
  ℹtem \textbf{Locks & Traps:} Mechanical = Mechanics + Attribute; Arcane = Arcana + Attribute.
  ℹtem \textbf{People vs. Mechanisms:} Subterfuge for social infiltration; Stealth to avoid detection; Mechanics or Arcana to bypass a physical or magical barrier.
  ℹtem \textbf{Formal vs. Informal:} Diplomacy for courts and treaties; Sway for bargaining and personal persuasion.
\end{itemize}
\end{tcolorbox}

\clearpage
\section*{Appendix C: Deck and Magic Reference}
\addcontentsline{toc}{section}{Appendix C: Deck and Magic Reference}
łabel{app:deck−magic}

\begin{tcolorbox}[title=Deck Types and Meanings, colback=white!97!gray, colframe=black!80!gray]
∈dex{decks}
\begin{description}
  ℹtem[Travel Decks (regional, 52−card)] Spade = Place, Heart = Actor, Club = Pressure, Diamond = Leverage.
  ℹtem[Deck of Consequences (scene drama)] Hearts = Social/Emotional, Spades = Harm/Escalation, Clubs = Material Cost, Diamonds = Magical/Spiritual.
\end{description}
\textbf{Important:} Never mix suit meanings across different deck types.
\end{tcolorbox}

\begin{tcolorbox}[title=Deck Usage Procedure, colback=white!97!gray, colframe=black!80!gray]
∈dex{deck of consequences!procedure}
After a roll that generates Story Beats:
\begin{enumerate}
  ℹtem \textbf{Direct Spend:} Translate SB into immediate consequences or timer ticks (see Appendix A).
  ℹtem \textbf{Optional Deck Draw:} Draw up to $\min(\text{SB},3)$ cards and synthesize a single narrative twist guided by the suits and the highest rank.
\end{enumerate}
\end{tcolorbox}

\begin{tcolorbox}[title=Rank Severity Guide, colback=white!97!gray, colframe=black!80!gray]
\begin{description}
  ℹtem[Ace−−3] Minor inconvenience or flavor.
  ℹtem[4−−6] Moderate setback.
  ℹtem[7−−9] Significant consequence.
  ℹtem[10−−King] Major fallout with lasting effects.
\end{description}
\end{tcolorbox}

\begin{longtable}{lp{10cm}}
∩tion{Deck of Consequences Interpretation} łabel{tab:consequences−interp} \\
⊤rule
\textbf{Suit} & \textbf{Interpretation} \\
∣rule
\endfirsthead
μlticolumn{2}{c}%
{{\tablename\ θble{} −− continued}} \\
⊤rule
\textbf{Suit} & \textbf{Interpretation} \\
∣rule
\endhead
⊥tomrule
\endfoot
Hearts (Emotional/Social) & Ace−−3: awkward social moment; 4−−6: relationship strain; 7−−9: betrayal; 10−−King: exile or shattered alliance. \\
Spades (Harm/Escalation) & Ace−−3: bruise or Fatigue; 4−−6: wound or damaged gear; 7−−9: severe injury; 10−−King: death or permanent loss. \\
Clubs (Material/Cost) & Ace−−3: minor resource loss; 4−−6: gear failure or debt; 7−−9: major asset lost; 10−−King: total ruin. \\
Diamonds (Magical/Spiritual) & Ace−−3: omen or whisper; 4−−6: curse or spirit appears; 7−−9: arcane backlash; 10−−King: reality bends temporarily. \\
\end{longtable}
∈dex{deck of consequences!interpretation}

\begin{longtable}{llll}
∩tion{Magic Paths Comparison} łabel{tab:magic−paths} \\
⊤rule
\textbf{Path} & \textbf{Requirements} & \textbf{Key Feature} & \textbf{Risk Type} \\
∣rule
\endfirsthead
μlticolumn{4}{c}%
{{\tablename\ θble{} −− continued}} \\
⊤rule
\textbf{Path} & \textbf{Requirements} & \textbf{Key Feature} & \textbf{Risk Type} \\
∣rule
\endhead
⊥tomrule
\endfoot
Caster (Freeform) & Spellcraft (6 XP) & Improvisational spellcasting & Backlash (\ref{subsec:elemental−backlash}) \\
Runekeeper (Rites) & Thiasos (2 XP) + Codex (4 XP) & Structured Patron Rites & Obligation (\ref{sec:obligation−corruption−gm}) \\
Invoker (Symbols) & Patron's Symbol (4 XP each) & Ritual precision, emergency casting & Symbol compromise (\ref{subsec:corruption−blooms}) \\
\end{longtable}
∈dex{magic!paths comparison}

\begin{tcolorbox}[title=DV Reference for Rites, colback=white!97!gray, colframe=black!80!gray]
∈dex{rites!difficulty value}
\[
\mathrm{DV} = \max(\text{ObligationCost} −− \text{Spirit},  \text{Tier})
\]

\begin{tabular}{c|cccc|c|c|c}
Obligation & Spirit 0 & Spirit 1 & Spirit 2 & Spirit 3−−4 & Tier 1 & Tier 2 & Tier 3 \\
ℎline
1 & 1 & 1 & 1 & 1 & 1 & 2 & 3 \\
2 & 2 & 1 & 1 & 1 & 1 & 2 & 3 \\
3 & 3 & 2 & 1 & 1 & 1 & 2 & 3 \\
4 & 4 & 3 & 2 & 1 & 1 & 2 & 3 \\
5 & 5 & 4 & 3 & 2 & 1 & 2 & 3 \\
6 & 6 & 5 & 4 & 3 & 1 & 2 & 3 \\
7 & 7 & 6 & 5 & 4 & 1 & 2 & 3 \\
\end{tabular}
\end{tcolorbox}

\begin{tcolorbox}[title=Casting Loop Summary, colback=white!97!gray, colframe=black!80!gray]
∈dex{free casting!procedure}
\begin{enumerate}
  ℹtem \textbf{Channel:} Wits + Arcana to gather magical potential.
  ℹtem \textbf{Weave:} Wits + Art (or appropriate magical skill) to shape the effect.
  ℹtem \textbf{Backlash:} The GM spends SB generated during the roll to impose thematic consequences (see \ref{subsec:elemental−backlash}).
\end{enumerate}
\end{tcolorbox}

\begin{tcolorbox}[title=Eight Elements of Magic, colback=white!97!gray, colframe=black!80!gray]
∈dex{elements of magic}
\begin{description}
  ℹtem[Earth] Solidity, stability, foundation.
  ℹtem[Fire] Energy, transformation, destruction.
  ℹtem[Air] Movement, speed, freedom.
  ℹtem[Water] Fluidity, healing, adaptability.
  ℹtem[Fate] Destiny, inevitability, causality.
  ℹtem[Life] Vitality, creation, growth.
  ℹtem[Luck] Chance, unpredictability.
  ℹtem[Death/Dreams] Endings, thresholds, the subconscious.
\end{description}
\end{tcolorbox}

\clearpage
\section*{Appendix D: Campaign Management}
\addcontentsline{toc}{section}{Appendix D: Campaign Management}
łabel{app:campaign}

\begin{tcolorbox}[title=Campaign Timer Examples, colback=white!97!gray, colframe=black!80!gray]
∈dex{campaign timers!examples}
\begin{itemize}
  ℹtem \textbf{Mandate Advancement:} Public victories, crisis resolution, recognition by powerful factions (see \ref{subsec:mandate}).
  ℹtem \textbf{Crisis Advancement:} Rival faction gains influence, asset neglect, betrayal, scandal (see \ref{subsec:crisis}).
\end{itemize}
\end{tcolorbox}

\begin{tcolorbox}[title=Follower and Asset Condition States, colback=white!97!gray, colframe=black!80!gray]
∈dex{follower conditions}∈dex{asset conditions}
\begin{description}
  ℹtem[Maintained] Full capability, no penalties.
  ℹtem[Neglected] −−1 die penalty; narrative wear and tear.
  ℹtem[Compromised] Unavailable until repaired or recovered through downtime.
\end{description}
See \ref{sec:followers−assets−gm} for detailed rules.
\end{tcolorbox}

\begin{tcolorbox}[title=Faction Loyalty Scale (Optional), colback=white!97!gray, colframe=black!80!gray]
∈dex{factions!loyalty scale}
Track faction standing on a simple −−3 to +3 scale.
\begin{description}
  ℹtem[−−3 Enemy] Actively works against the party.
  ℹtem[0 Neutral] Indifferent.
  ℹtem[+3 Ally] Will sacrifice for the party.
\end{description}
Adjust by ±1 after major interactions (saving or betraying the faction).
\end{tcolorbox}

\begin{tcolorbox}[title=Campaign Pacing (Arc Structure), colback=white!97!gray, colframe=black!80!gray]
∈dex{campaign pacing}
\begin{enumerate}
  ℹtem \textbf{Introduction} (1−−2 sessions): Establish stakes, hook the players.
  ℹtem \textbf{Development} (2−−4 sessions): Complications multiply, alliances form.
  ℹtem \textbf{Climax} (1−−2 sessions): Major confrontation, critical choices.
  ℹtem \textbf{Resolution} (1 session): Consequences land, new status quo established.
\end{enumerate}
\end{tcolorbox}

\clearpage
\section*{Appendix E: Safety and Inclusivity}
\addcontentsline{toc}{section}{Appendix E: Safety and Inclusivity}
łabel{app:safety}

\begin{tcolorbox}[title=X−Card System, colback=white!97!gray, colframe=black!80!gray]
∈dex{safety tools!X−Card}
\begin{itemize}
  ℹtem \textbf{X−Card:} ``I'm not comfortable −− change direction immediately.'' No explanation required.
  ℹtem \textbf{O−Card:} ``I'm excited −− lean into this element.''
  ℹtem \textbf{N−Card:} ``I need a break −− pause the scene.''
\end{itemize}
Keep cards visible. Model appropriate use as the GM.
\end{tcolorbox}

\begin{tcolorbox}[title=Script Change, colback=white!97!gray, colframe=black!80!gray]
∈dex{safety tools!script change}
Ask at appropriate moments:
\begin{itemize}
  ℹtem \textbf{Pause:} ``Are you okay with this content?''
  ℹtem \textbf{Rewind:} ``Would you like to go back and change something?''
  ℹtem \textbf{Fast Forward:} ``Would you like to skip ahead?''
\end{itemize}
\end{tcolorbox}

\begin{tcolorbox}[title=Inclusive Worldbuilding, colback=white!97!gray, colframe=black!80!gray]
∈dex{inclusivity}
\begin{itemize}
  ℹtem Research real cultures thoroughly when drawing inspiration.
  ℹtem Avoid stereotypes; build characters with depth and agency.
  ℹtem Provide character options that do not rely on real−world cultural markers.
  ℹtem Use high−contrast, readable materials for accessibility.
\end{itemize}
\end{tcolorbox}

\begin{tcolorbox}[title=Conflict Resolution, colback=white!97!gray, colframe=black!80!gray]
∈dex{conflict resolution}
When conflicts arise at the table:
\begin{enumerate}
  ℹtem Pause gameplay.
  ℹtem Discuss privately if needed.
  ℹtem Focus on understanding, not assigning fault.
  ℹtem Find collaborative solutions.
  ℹtem Document any agreements.
\end{enumerate}
Common solutions: the GM's ruling stands for the session; research the rule later; adjust future play.
\end{tcolorbox}

\clearpage
\section*{Appendix F: Rival Party Generator}
\addcontentsline{toc}{section}{Appendix F: Rival Party Generator}
łabel{app:rival−generator}
∈dex{rival parties!generator}

Use the following d6 tables to generate a rival party in under 5 minutes. Roll once in each category or choose deliberately.

\begin{tcolorbox}[title=Rival Party Name, colback=white!97!gray, colframe=black!80!gray]
¢ering
\setlength{\tabcolsep}{6pt}
\begin{tabular}{cll}
⊤rule
\textbf{d6} & \textbf{Name} & \textbf{Flavor} \\
∣rule
1 & The Ashen Cord & Echoes, smoke, unbreakable knots. \\
2 & The Brass Legion & Gears, discipline, cold precision. \\
3 & The Thorns of Winter & Ice, silence, patient cruelty. \\
4 & The Gilded Chain & Wealth, obligation, gilded cages. \\
5 & The Salt Pact & Oaths sworn in brine, pirates or pilgrims. \\
6 & The Unfaced & Masks, secrets, stolen identities. \\
⊥tomrule
\end{tabular}
\end{tcolorbox}

\begin{tcolorbox}[title=Composition (3−−5 members), colback=white!97!gray, colframe=black!80!gray]
¢ering
\begin{tabular}{cll}
⊤rule
\textbf{d6} & \textbf{Role} & \textbf{Archetype} \\
∣rule
1 & Leader & Shadow mage, war−priest, strategist, or oath−breaker. \\
2 & Muscle & Berserker, knight, duelist, or beast−handler. \\
3 & Mage/Weird & Cantor, free caster, runekeeper, or witch. \\
4 & Scout/Thief & Infiltrator, poisoner, sniper, or ghost. \\
5 & Face/Healer & Negotiator, spy, medic, or confessor. \\
6 & Specialist & Engineer, summoner, trapper, or lore−master. \\
⊥tomrule
\end{tabular}
\end{tcolorbox}

\begin{tcolorbox}[title=Patron (roll once per character or once for the party), colback=white!97!gray, colframe=black!80!gray]
¢ering
\begin{tabular}{cll}
⊤rule
\textbf{d6} & \textbf{Patron} & \textbf{Theme} \\
∣rule
1 & Ikasha (Shadow) & Secrets, thresholds, omens. \\
2 & Malachai (Curse) & Corrupted bargains, venomous gifts. \\
3 & The Traveler (Roads) & Journeys, shortcuts, debts. \\
4 & The Ninth (Knowledge) & Forbidden truths, fractured minds. \\
5 & Khemesh (Depths) & Unstoppable pressure, drowning. \\
6 & Palinode (Performance) & Rapture, encores, hungry fame. \\
⊥tomrule
\end{tabular}
\end{tcolorbox}

\begin{tcolorbox}[title=Goal (pressure or obstacle), colback=white!97!gray, colframe=black!80!gray]
¢ering
\begin{tabular}{cll}
⊤rule
\textbf{d6} & \textbf{Goal} & \textbf{Example} \\
∣rule
1 & Retrieve the same artifact. & Sunstone, charter, relic. \\
2 & Capture/kill a mutual contact. & Witness, defector, rival. \\
3 & Secure a Patron's favor. & Complete a ritual or sacrifice. \\
4 & Eliminate the PCs for a rival. & Bounty, vendetta, jealous patron. \\
5 & Claim a contested site. & Anchorhold, shrine, trade route. \\
6 & Prevent a prophecy. & Stop the PCs from changing fate. \\
⊥tomrule
\end{tabular}
\end{tcolorbox}

\begin{tcolorbox}[title=Relationship to PCs, colback=white!97!gray, colframe=black!80!gray]
¢ering
\begin{tabular}{cll}
⊤rule
\textbf{d6} & \textbf{Relationship} & \textbf{Initial Attitude} \\
∣rule
1 & Unexpected rivals & Neutral −− will negotiate first. \\
2 & Hunters seeking a bounty & Hostile −− will attack if confident. \\
3 & Former allies, now betrayed & Bitter −− may accept apology at cost. \\
4 & Pawns of a common enemy & Reluctant allies −− temporary truce. \\
5 & Zealous competitors & Competitive −− will cheat but not kill. \\
6 & Mysterious & Unreadable −− may help or hinder. \\
⊥tomrule
\end{tabular}
\end{tcolorbox}

\begin{tcolorbox}[title=Additional Resource (off−screen), colback=white!97!gray, colframe=black!80!gray]
¢ering
\begin{tabular}{cll}
⊤rule
\textbf{d6} & \textbf{Resource} & \textbf{Effect} \\
∣rule
1 & Hidden safehouse & Can hide from pursuers once. \\
2 & Bribed official & May delay a legal timer once. \\
3 & Spy network & Gains +1 die to gather info. \\
4 & War chest & Can hire temporary muscle. \\
5 & Patron's omen & Once per arc, avoid a Miss. \\
6 & Stolen map & Shortcut to any location once. \\
⊥tomrule
\end{tabular}
\end{tcolorbox}

\subsection*{Rival Party Timer }
Every rival party should have a visible agenda timer, e.g., \emph{Rival Agenda [6]}. Tick it:
\begin{itemize}
  ℹtem When the party ignores or fails to hinder the rivals.
  ℹtem When the GM spends 2+ SB on off−screen rival activity.
  ℹtem When a session ends without contact.
\end{itemize}
When the timer fills, the rivals achieve a major milestone (steal the relic, secure the Patron's favor, etc.). The GM reveals their progress through rumors, looted sites, or a direct confrontation.

\subsection*{Example Rival Party (Quick Build)}
\begin{quote}
\textbf{The Ashen Cord} (from simulation)\\
\emph{Composition:} Vesper (leader, shadow mage, serves Ikasha), Tallow (berserker, serves Malachai), Ember (cantor, serves Palinode), Cinder (scout, no Patron).\\
\emph{Goal:} Retrieve the Sunstone fragment before the PCs.\\
\emph{Relationship:} Unexpected rivals −− will negotiate or betray as needed.\\
\emph{Clock:} \emph{Cord's Agenda [6]}.\\
\emph{Resource:} Hidden safehouse in Silkstrand.
\end{quote}

\clearpage
\section*{Appendix G: One−Page GM Screen Summary}
\addcontentsline{toc}{section}{Appendix G: One−Page GM Screen Summary}
łabel{app:gm−screen}

\begin{tcolorbox}[title=\textbf{Fate's Edge −− GM Screen}, colback=blue!5!white, colframe=blue!75!black, sharp corners, fonttitle=\bfseries, center title]
\begin{center}
łarge\textbf{CORE RESOLUTION}
\end{center}
\begin{tabularx}{\textwidth}{X|X}
ℎline
\textbf{1. Declare Action} & Intent + Attribute + Skill \\
\textbf{2. Set DV} & 2=Routine, 3=Pressured, 4=Hard, 5+=Extreme \\
\textbf{3. Set Position} & Dominant (re−roll failure), Controlled (normal), Desperate (re−roll success) \\
\textbf{4. Roll} & Count successes (6+) and 1s (SB for GM) \\
\textbf{5. Outcome} & Clean Success, Success & Cost, Partial (1 Boon), Miss (2 Boons) \\
\textbf{6. Spend SB} & 1=minor, 2=moderate, 3+=major complication \\
ℎline
\end{tabularx}

\vspace{0.5em}
\begin{center}
łarge\textbf{COMBAT QUICK REF}
\end{center}
\begin{tabularx}{\textwidth}{X|X}
ℎline
\textbf{Move} & Shift one band (Close\(\(↔\)\)Near\(\(↔\)\)Far) \\
\textbf{Dash} & Move two bands, next defense Desperate \\
\textbf{Charge/Skirmish} & Move two bands + attack, cost 1 Fatigue, requires open ground \\
\textbf{Opportunity Attack} & Leaving Close without Disengage = free melee attack \\
\textbf{Fatigue} & Each Fatigue worsens Position; fill → Harm +1, clear Fatigue \\
\textbf{Harm} & 1=−−1 die (major), 2=−−1 dice (most), 3=incapacitated \\
ℎline
\end{tabularx}

\vspace{0.5em}
\begin{center}
łarge\textbf{MAGIC SHORTCUTS}
\end{center}
\begin{tabularx}{\textwidth}{X|X}
ℎline
\textbf{Free Caster} & DV = 1 + number of [TAGS] (dangerous +2); backlash on 1s \\
\textbf{Runekeeper} & Invoke Rite (1 action) +1 Obligation; Push It: +1 Obligation \\
\textbf{Invoker} & Ritual (DV+1 rounds) +1 Obligation; Crack the Seal: instant, +2 Obligation, Symbol Compromised \\
\textbf{Cantor} & Song (1 action, resolves next turn) +Corruption timer; Push: resolve immediately \\
\textbf{Summoning} & Call (1 action) + Bind (1 Boon or Fatigue); Leash = Cap + Spirit segments \\
ℎline
\end{tabularx}

\vspace{0.5em}
\begin{center}
łarge\textbf{SB SPEND QUICK}
\end{center}
\begin{tabularx}{\textwidth}{X|X}
ℎline
\textbf{1 SB} & Noise, trace, tick a timer +1 \\
\textbf{2 SB} & Alarm raised, lose advantage, lesser foe appears \\
\textbf{3 SB} & Reinforcements, key gear break, campaign front tick \\
\textbf{4+ SB} & Trap springs, authority arrives, scene shift \\
ℎline
\end{tabularx}

\vspace{0.5em}
\begin{center}
łarge\textbf{CLOCKS & DIFFICULTY}
\end{center}
\begin{tabularx}{\textwidth}{X|X}
ℎline
\textbf{Timer Size} & 4=short, 6=standard, 8=long, 10=epic \\
\textbf{Default DV} & 3 (pressured, mild opposition) \\
\textbf{DV 2} & Routine, trivial stakes \\
\textbf{DV 4} & Hard, active opposition \\
\textbf{DV 5+} & Extreme, high drama \\
ℎline
\end{tabularx}

\vspace{0.5em}
\small¢ering \textit{Golden Rule: When in doubt, make a ruling that keeps the story moving.}
\end{tcolorbox}

\clearpage
\section*{Appendix H: Sample Deck Reading (Example)}
\addcontentsline{toc}{section}{Appendix H: Sample Deck Reading (Example)}
łabel{app:sample−reading}
∈dex{deck of consequences!sample reading}

The GM draws three cards from the Deck of Consequences after a player rolls three 1s on a critical stealth check.

\textbf{Drawn:} Knave of Hearts, 7 of Clubs, Ace of Spades.

\paragraph{Synthesis:} The GM combines the suits:
\begin{itemize}
  ℹtem \textbf{Hearts (social)} + \textbf{Spades (harm)} = betrayal with injury.
  ℹtem \textbf{Clubs (resource cost)} = material loss.
  ℹtem \textbf{Ace} = scene−altering twist.
\end{itemize}

\paragraph{GM narration:}
\begin{quote}
\textquotedbl You slip past the first guard, but your cloak catches on a rusty nail (Clubs −− resource drain: your cloak tears, leaving a scrap of fabric behind). Worse, the second guard recognizes your face from a wanted poster −− she's your estranged cousin (Hearts −− social betrayal). She hesitates for only a heartbeat, then shouts, ``Traitor!'' (Spades −− harm escalates). As she draws her blade, a hidden portcullis drops behind you, sealing off the exit (Ace −− scene shift). You're trapped in the courtyard with her and two more guards.\textquotedbl
\end{quote}

\noindent This single three−card draw creates a layered, memorable complication from a single SB spend.

\clearpage
\section*{Appendix I: Session Zero Checklist Handout}
\addcontentsline{toc}{section}{Appendix I: Session Zero Checklist Handout}
łabel{app:session−zero−handout}
∈dex{Session Zero!handout}

\begin{tcolorbox}[colback=green!5!white, colframe=green!75!black, title=Session Zero Checklist, fonttitle=\bfseries]
Use this checklist during your Session Zero. Write down the answers where everyone can see them.
\begin{enumerate}
  ℹtem \textbf{Lines & Veils} −− What content is off−limits (Lines) or only implied (Veils)?
  ℹtem \textbf{Tone} −− Cinematic action? Grim survival? Political intrigue? Dark horror? Mix of several?
  ℹtem \textbf{Campaign length} −− One−shot, short arc (3−−6 sessions), or long campaign (12+ sessions)?
  ℹtem \textbf{Character hooks} −− Each player shares one goal and one fear for their character.
  ℹtem \textbf{Patron interest} −− Which Patrons (if any) are already watching the party?
  ℹtem \textbf{Starting situation} −− Where does the story begin? (Use the Crown Spread or a simple opener.)
  ℹtem \textbf{House rules} −− Any optional subsystems (Threat Pool, Deck of Consequences, etc.) you plan to use?
  ℹtem \textbf{Scheduling & attendance} −− How often do you meet? What happens if a player misses a session?
  ℹtem \textbf{Player−GM contract} −− Remind the table: \textquotedbl We are co−creators, not adversaries.\textquotedbl
  ℹtem \textbf{Resource mindset} −− \textquotedbl Boons are not XP tokens. Spend them. Embrace failure.\textquotedbl
  ℹtem \textbf{Patron expectation} −− \textquotedbl They are allies who demand payment −− embrace it.\textquotedbl
  ℹtem \textbf{Flavor is free} −− Encourage descriptive narration for all actions, not just magic.
\end{enumerate}
\end{tcolorbox}

% =========================================================
% END OF APPENDICES (CHAPTER 12)
% =========================================================

\clearpage
\chapter{Using Expansions: A Guide for the Overwhelmed GM}
łabel{ch:using−expansions}
∈dex{expansions}

\begin{quote}
\small
\textit{\textquotedbl I bought every expansion. Then I sat at my table with thirty−seven PDFs and a blank look. The players were waiting. The beer was warm. I learned that night: you don't use expansions. You borrow from them.\textquotedbl}
−−− The Grey Wanderer, from a letter found under a stack of printouts
\end{quote}

\section{The First Rule: You Own the Expansions, They Don't Own You}

Every expansion in Fate's Edge is a toolbox, not a syllabus. You are never required to use a single rule, timer, or cultural detail from any expansion. The core game runs perfectly well on its own. Expansions exist to \emph{solve specific problems} or \emph{add specific flavors} when you need them.

\subsection{The Three Questions to Ask Before Opening Any Expansion}

Before you reach for a PDF, ask yourself:

\begin{enumerate}
ℹtem \textbf{What problem am I solving?} (e.g., ``combat feels same−y'', ``my players want more political intrigue'', ``we're going to sea'')
ℹtem \textbf{What is the smallest piece of this expansion that solves it?} (e.g., ``just the Seven Bell Court judging table'', ``just the travel timers'', ``just the Lethai−ar witness knots'')
ℹtem \textbf{Can I explain that piece to my players in one minute?} If not, pick a smaller piece.
\end{enumerate}

\subsection{The 10% Rule}

A typical Fate's Edge expansion is 30−100 pages. In a single session, you will use \textbf{at most 10% of any expansion}. The other 90% is background, lore, alternative systems, and optional depth. You are not expected to master all of it.

\section{A Taxonomy of Expansions (What They're Actually For)}

Use this table to quickly identify which expansion to reach for. Expansions are grouped by \textbf{purpose}, not by title.

\begin{longtable}{p{3.5cm}|p{3.5cm}|p{6cm}}
ℎline
\textbf{If you want...} & \textbf{Use expansions in this family} & \textbf{Example of a single mechanic to steal} \\
ℎline
Honorable duels, tournaments, social combat & Seven Bell Court, Vhasia: The Knave's Regret & The Form/Spirit/Intent judging table (3 rows, no dice) \\
ℎline
New martial options (no magic) & The Way of the Warrior, Ykrul: Iron & Blood & A single talent chain (e.g., Sentinel, Bladestorm, Pathfinder) \\
ℎline
Cultural depth for a specific people & Eastern Realms, Peoples and Cultures, The Frost Ledger & One etiquette hook and one String (e.g., ``Aelinnel never ring the ninth bell'') \\
ℎline
Investigation & The Salt Ledger, The Wandering Ledger, Saikou Ira's Compendium & One clue type (e.g., ``a silver coin in the mouth = ferryman's toll'') \\
ℎline
Sea travel, naval combat, port intrigue & The Amaranthine Sea & Travel timers [4−8] and the Inspection procedure \\
ℎline
Folk horror, psychological tension & Horror Campaigns, The Frost Ledger, Mistslands expansion & The Dread Timer (4−8 segments) \\
ℎline
Witchcraft, curses, bargains & Witches of Fate's Edge, Daughters of the Unbroken Cord & One working type (e.g., Chord Working, Promise Timer , Cord of Flesh) \\
ℎline
Patron−heavy campaigns & The Grey Wanderer's Grimoire, Malachai, The Chained Angel & One new Patron with its Rites and Obligation flavor \\
ℎline
Dungeon crawling, ancient lairs & Dragon's Lair, The Book of Air, Mist, Alder, Thorn, and Mirror & Lair moves (3−4 per scene) \\
ℎline
Political intrigue & Political Intrigue framework, Vhasia, Viterra expansions & Faction timers (Influence/Stability/Exposure) \\
ℎline
Caravans, travel, trade & Caravans: The Way of Silk, The Wandering Ledger & Travel intent dial (Speed/Stealth/Show/Tribute) \\
ℎline
\end{longtable}

\section{The Hybrid Adventure: How to Cook with Leftovers}

The most successful way to use expansions is to build a \textbf{hybrid adventure}: one core story (often from a free adventure or your own imagination) with mechanics borrowed from 2−4 expansions.

\subsection{Example: The Tournament of the Broken Crown}

The GM wanted:
\begin{itemize}
ℹtem A tournament (so players could fight without dying).
ℹtem A murder mystery (so non−combat characters had something to do).
ℹtem Cultural flavor (so the world felt lived−in).
\end{itemize}

They borrowed:
\begin{itemize}
ℹtem \textbf{Judging table} from \textit{Seven Bell Court} (3 rows: Form, Spirit, Intent).
ℹtem \textbf{Unbuckled gorget clue} from \textit{Vhasia: The Knave's Regret} (mercy as a theme).
ℹtem \textbf{Bell−tending} and \textbf{witness knots} from \textit{Eastern Realms} (investigation hooks).
ℹtem \textbf{Martial talents} from \textit{The Way of the Warrior} (Disciplined Body, Shield Bearer, Flurry).
\end{itemize}

\textbf{Prep time:} 90 minutes. \textbf{Pages read:} ~15 total across four expansions. \textbf{Result:} A 3−hour session that felt fresh, deep, and manageable.

\subsection{The Hybrid Recipe}

\begin{enumerate}
ℹtem Pick a \textbf{core activity} (tournament, heist, travel, investigation).
ℹtem Identify \textbf{one gap} in the core rules that an expansion fills (e.g., ``combat lacks stakes beyond HP'' → Seven Bell Court).
ℹtem Read \textbf{only the first 5−10 pages} of that expansion until you find a mechanic you understand.
ℹtem Add \textbf{one cultural hook} from a lore expansion (e.g., ``in this city, knights unbuckle their gorgets to show mercy'').
ℹtem Add \textbf{one talent} from a player−options expansion (e.g., give an NPC a talent from \textit{The Way of the Warrior}).
ℹtem Run the session. \textbf{Ignore every other rule} from the expansions.
\end{enumerate}

\section{Managing Player Expectations}

When players know you own expansions, they will ask to use things. This is a gift, not a burden.

\subsection{The ``Yes, But Later'' Rule}

When a player asks ``Can I use the Shadow Step talent from \textit{Witches of Fate's Edge}?''

\begin{itemize}
ℹtem \textbf{Yes:} ``Yes, but you'll need to earn it through roleplay. Find a hag or a Lethai−ar teacher in the next town.''
ℹtem \textbf{No, but:} ``No, but you can take a different talent from the core book that works similarly, and we'll reskin it.''
ℹtem \textbf{Later:} ``Not this session, but tell me about it during downtime and we'll build it into your character's arc.''
\end{itemize}

\subsection{The Expansion Session Zero}

At the start of a new campaign or arc, ask your players:

\begin{enumerate}
ℹtem ``Which expansions have you read?'' (Be honest; they may have read none.)
ℹtem ``What's one thing from an expansion you'd love to see in play?'' (Write it down.)
ℹtem ``What's one thing you absolutely don't want?'' (Lines and veils for expansion content.)
\end{enumerate}

Then pick \textbf{one expansion} to feature prominently in the next 2−3 sessions. Rotate.

\section{Cross−Expansion Synergy (When Two Books Work Better Than One)}

Some expansions are designed to work together. Others accidentally complement each other. Here are known synergies:

\begin{itemize}
ℹtem \textbf{Seven Bell Court + Vhasia: The Knave's Regret} = duels that are judged on honor, not just damage. The ``unbuckled gorget'' becomes a mechanical choice (lower Defense, higher Spirit points).
ℹtem \textbf{The Amaranthine Sea + The Salt Ledger} = sea travel with port mysteries. Use the travel timers from \textit{Sea} and the lighthouse/spirit stories from \textit{Salt} as random encounters.
ℹtem \textbf{Witches of Fate's Edge + Horror Campaigns} = curse−based horror. Use the Dread Timer and the Cord mechanics together; each curse failure ticks the Dread Timer .
ℹtem \textbf{Eastern Realms + The Way of the Warrior} = culturally−infused martial characters. Let a Lethai−al samurai take Bushido−style talents from \textit{Warrior} and reflavor them as ``the Way of the Leaf.''
\end{itemize}

\subsection{Antisynergy: When Not to Mix}

Some expansions actively fight each other. Avoid:

\begin{itemize}
ℹtem \textbf{Too many timers.} \textit{Horror Campaigns} (Dread Timer , Reality Fracture) + \textit{Political Intrigue} (Mandate, Crisis) + \textit{The Amaranthine Sea} (Weather Window, Patrol Net) = 6+ timers. Players cannot track that many. Pick 2−3 maximum per session.
ℹtem \textbf{Contradictory cultural lore.} The \textit{Eastern Realms} and \textit{The Salt Ledger} have different versions of Aelinnel bell−customs. Choose one interpretation for your campaign and stick to it. Tell players ``we're using the \textit{Eastern Realms} version.''
\end{itemize}

\section{Practical Tools: The Expansion Integration Checklist}

Before a session where you plan to use expansion content, run through this checklist.

\subsection{Pre−Session (30−60 minutes)}

\begin{enumerate}
ℹtem \textbf{Select 1−3 expansions} to draw from. Write their names on an index card.
ℹtem \textbf{Extract 1−3 mechanics} from each expansion. Write each mechanic as a single sentence on the card.
ℹtem \textbf{Prepare one page of reference.} Print the relevant table (e.g., Seven Bell Court judging, travel timer sizes, talent list). Do not print the whole expansion.
ℹtem \textbf{Note page numbers.} If you must look something up, know where it is.
ℹtem \textbf{Check for conflicts.} Do any of your chosen mechanics use the same resource (e.g., both use Boons in contradictory ways)? If yes, pick one.
\end{enumerate}

\subsection{During Session (At the Table)}

\begin{enumerate}
ℹtem \textbf{Announce the new mechanic.} ``Tonight we're using the Seven Bell Court. Your duels will be judged on Form, Spirit, and Intent.''
ℹtem \textbf{Demonstrate once.} Run a quick mock duel or skill challenge to show how it works.
ℹtem \textbf{Let players opt out.} If someone doesn't want to use a new subsystem, let them use the core rules instead.
ℹtem \textbf{Keep a timer.} If a new mechanic slows play for more than 2 minutes, simplify it or drop it for the scene.
\end{enumerate}

\subsection{Post−Session (5 minutes)}

\begin{enumerate}
ℹtem \textbf{Ask:} ``What expansion mechanic worked? What didn't?''
ℹtem \textbf{Note:} Write down which pages you actually used. Toss the rest.
ℹtem \textbf{Decide:} Will you use this expansion again? If yes, put the relevant page in your GM binder. If no, archive it.
\end{enumerate}

\section{DM's Cheat Sheet: Quick Reference by Page Count}

If you have limited time before a session, use this triage guide.

\begin{longtable}{p{4cm}|p{4cm}|p{4cm}}
ℎline
\textbf{Expansion} & \textbf{Read these pages} & \textbf{Skip these pages} \\
ℎline
Seven Bell Court & 5−7 (judging table, Honor/Pragmatism) & Tournament logistics, NPC champion statblocks \\
ℎline
Vhasia: The Knave's Regret & 1−3 (the opening fiction), 25−27 (talents) & Full canton descriptions, duel history \\
ℎline
Way of the Warrior & 31 (talent quick reference) & Cultural warrior variants (unless a player is one) \\
ℎline
Eastern Realms & 13−14 (bell−tending, witness knots) & Travel decks (use only when traveling) \\
ℎline
Witches of Fate's Edge & i−iv (first principles), 89−91 (Threshold working) & Deep hag lore, full talent trees \\
ℎline
The Amaranthine Sea & 9−11 (travel timers, inspection) & Fleet combat rules, ship statblocks \\
ℎline
Horror Campaigns & 5−6 (Dread Timer ), 8−9 (breaking points) & Extended campaign timers (use only for long arcs) \\
ℎline
Dragon's Lair & 3−5 (lair generator), 21 (dragon moves) & Hoard generators, dragon stats (improvise Tier) \\
ℎline
The Salt Ledger & Any single story (use as adventure seed) & The rest of the stories (save for later) \\
ℎline
\end{longtable}

\section*{Magic Adjudication Quick Reference (GM)}
łabel{sec:gm−magic−cheat}

\begin{tcolorbox}[colback=red!5!white,colframe=red!75!black,title=Default Magic DV]
¢ering
\begin{tabular}{ll}
\textbf{Effect scope} & \textbf{DV} \\
Cantrip (1 tag) & 2 \\
Combat spell (2 tags) & 3 \\
Ritual (small, 3 tags) & 4 \\
Ritual (large, dangerous tags) & 5+ \\
\end{tabular}
\end{tcolorbox}

\subsection*{Free Casting (TAGS)}
\[
DV = 1 + \text{#tags}  (\text{dangerous tags: } +2\text{ instead of }+1)
\]
\emph{Dangerous tags:} Leap, Fold, Gate, Gravity, Create, Summon, Transmute, Animate, Transform, Dominate, Rewind, Resurrect.

\paragraph{Backlash:} Spend SB from 1s. Severity by Position:
\begin{itemize}
  ℹtem Dominant: minor (noise, Fatigue 1)
  ℹtem Controlled: moderate (Harm 1, Condition)
  ℹtem Desperate: severe (Harm 2, scene shift)
\end{itemize}

\subsection*{Rites (Runekeepers & Invokers)}
\[
DV = \max(\text{ObligationCost} − \text{Spirit},  \text{Tier: Low=1, Standard=2, High=3})
\]
\paragraph{Obligation overflow:} Exceed Capacity (Spirit+Presence) → mark 1 Fatigue per segment. Exceed 2× Capacity → +1 Harm, Patron intrusion.

\subsection*{Healing Magic}
\begin{itemize}
  ℹtem \textbf{[HEAL] only:} Restores 1 Fatigue. DV 2. 1 action.
  ℹtem \textbf{[HEAL] + another tag (Purify, Strengthen, Cleanse):} Restores 1 Harm. DV = 2 + number of tags, plus \textbf{+2 surcharge}. Requires two actions unless Push.
  ℹtem \textbf{Ritual healing (Harm 2+):} Requires Extended ritual (timer), rare components.
\end{itemize}

\subsection*{Summoning}
\begin{itemize}
  ℹtem \textbf{Call:} 1 action, spirit manifests.
  ℹtem \textbf{Bind:} 1 Boon or 1 Fatigue.
  ℹtem \textbf{Leash capacity:} Cap + Spirit segments. Tick when: spirit takes Harm, acts against nature, crosses a [WARD], summoner splits focus, rival contests control.
  ℹtem \textbf{Boon Finesse:} Once/round, spend 1 Boon to clear 1 Leash tick.
  ℹtem \textbf{Leash full:} Spirit acts once to its nature, then departs (or turns hostile).
\end{itemize}

\subsection*{Cantors (Songs)}
\begin{itemize}
  ℹtem \textbf{Sing:} 1 action, resolves next turn. Roll Lore + Performance (DV 2−−3).
  ℹtem \textbf{Push:} Resolve immediately, mark 1 Fatigue, advance Corruption Timer by 1.
  ℹtem \textbf{Corruption Timer :} Size = Spirit. When full → bloom (gain benefit & burden from last patron), reset to Tier (min 1).
  ℹtem \textbf{Resonant Rite:} Performance with 10+ witnesses → may found Chorus (minor follower benefit, Obligation to patron each session).
\end{itemize}

\subsection*{Patron Intrusions (when Obligation overflows or GM spends 2+ SB)}
\begin{itemize}
  ℹtem Demand a quest (``Find my stolen symbol.'')
  ℹtem Impose a cost (``Your next Rite costs double.'')
  ℹtem Send an omen (``Candles burn blue; you feel watched.'')
  ℹtem Summon a rival servant (``Another follower of Malachai is hunting you.'')
\end{itemize}

\subsection*{Common Magic Rulings}
\begin{itemize}
  ℹtem \textbf{Countermagic:} [COUNTER] tag vs. DV = spell's DV. Success cancels.
  ℹtem \textbf{Dispel:} [DISPEL] vs. DV = effect's Tier (1−−5). On hit, effect ends.
  ℹtem \textbf{Wards:} [WARD] vs. Outsider. DV = target's Cap. Crosser adds DV to Leash.
  ℹtem \textbf{Line of sight:} Required for most spells; without it +1 DV.
\end{itemize}

\section{Adversary Templates}

\subsection{Generic Adversaries}
łabel{subsec:generic−adversaries}
Use these for filler encounters. No special traits−−−just core stats. Scale by adjusting Threat Level (see Table~\ref{tab:threat−scaling}). All use standard attack/defense as per Chapter~\ref{ch:gm−core}.

\begin{itemize}
    ℹtem \textbf{Bandit Rabble} (Threat Level 1) \\
    \textit{Stats:} Body 2, Wits 1, Spirit 1, Presence 1   |   \textit{Skills:} Melee 1, Athletics 1 \\
    \textit{Flavor:} ``Desperate locals with rusty blades. They flee if half their number falls.'' \\
    \textit{Use:} Random encounters, tavern brawls, distractions.

    ℹtem \textbf{City Watch Recruit} (Threat Level 2) \\
    \textit{Stats:} Body 2, Wits 1, Spirit 1, Presence 2   |   \textit{Skills:} Melee 1, Command 1 \\
    \textit{Flavor:} ``Poorly trained but eager to prove themselves. Will call for backup if overwhelmed (alarm raises a 4−segment timer).'' \\
    \textit{Use:} Urban patrols, gate guards, low−level enforcement.

    ℹtem \textbf{Cult Novice} (Threat Level 1) \\
    \textit{Stats:} Body 1, Wits 1, Spirit 2, Presence 1   |   \textit{Skills:} Lore 1 (cult−specific), Sway 1 \\
    \textit{Flavor:} ``True believer, poorly trained. Will gladly die for the cause but breaks if shown proof their faith is a lie.'' \\
    \textit{Use:} Guarding rituals, spreading propaganda, cannon fodder. On a Miss when trying to recruit or intimidate, they may flee or surrender (GM discretion).

    ℹtem \textbf{Street Thug} (Threat Level 2) \\
    \textit{Stats:} Body 3, Wits 1, Spirit 1, Presence 1   |   \textit{Skills:} Brawl 2, Intimidation (Presence+Brawl) \\
    \textit{Flavor:} ``Enforcer for a local gang. Favors cheap shots and overwhelming numbers.'' \\
    \textit{Use:} Alley ambushes, protection rackets, gang warfare.

    ℹtem \textbf{Guard Sergeant} (Threat Level 3) \\
    \textit{Stats:} Body 3, Wits 2, Spirit 2, Presence 2   |   \textit{Skills:} Melee 2, Command 1, Athletics 1 \\
    \textit{Flavor:} ``Veteran who has seen it all. Leads from the front and shouts orders.'' \\
    \textit{Use:} Leader of a watch patrol, military outpost, or noble's bodyguard.
\end{itemize}

\begin{table}[htbp]
¢ering
∩tion{Threat Level Scaling}
łabel{tab:threat−scaling}
\begin{tabular}{lcll}
⊤rule
\textbf{Threat Level} & \textbf{Stat Adjustment} & \textbf{Typical Use} & \textbf{PC Tier Match} \\
∣rule
0 & −−1 to all stats & Non−threatening flavor & Below PC tier \\
1 & Base stats (as written) & Standard minions & PC tier −−1 \\
2 & +1 to all stats & Tougher minions / elite & PC tier \\
3 & +2 to all stats & Mini−boss equivalent & PC tier +1 \\
4+ & +3 to all stats & add one archetype trait & Major boss & PC tier +2 \\
⊥tomrule
\end{tabular}
\end{table}
\par\smallskip\textit{Adjust all Attributes and Skills uniformly. For Threat Level 0, treat as flavor only (no combat threat). For Threat Level 3+, consider giving one simple trait from the archetypes below.}

\subsection{Archetypal Adversaries}
łabel{subsec:archetypal−adversaries}
Each has one thematic mechanic tied to existing rules. Use as recurring threats, scene−setters, or mini−bosses. Scale via Threat Level table.

⫕section{The Iron Patriarch / Matriarch}
\textit{Archetype:} Ruthless leader of a family, guild, or corporation. Values legacy and control above all.\\
\textit{Stats (Threat Level 2):} Body 2, Wits 3, Spirit 2, Presence 3   |   \textit{Skills:} Sway 2, Command 1, Lore 1 (family/guild)\\
\textbf{Thematic Mechanic:} \textit{Legacy's Weight} \\
When the Patriarch/Matriarch would take Harm, they may immediately order a loyal subordinate (if present) to take the blow. Roll Wits + Command \textit{vs. DV 2}. On success, the subordinate takes 1 Harm instead (armor applies normally). On failure, the Patriarch takes the Harm but gains +1 die on their next action (resolve fueled by desperation).\\
\textit{Flavor Hook:} ``They don't fear death−−−they fear their life's work dying with them. Every blow you strike feels like it's aimed not at them, but at the empire they bled to build.''\\
\textit{Rules Tie−In:} Uses opposed rolls (see Chapter~\ref{ch:gm−core}) and Harm (see Chapter~\ref{ch:gm−core}).

⫕section{The Hollow Prophet}
\textit{Archetype:} Charismatic leader of a movement built on a lie (or partial truth). Believes their own deception.\\
\textit{Stats (Threat Level 2):} Body 1, Wits 2, Spirit 3, Presence 3   |   \textit{Skills:} Sway 2, Lore 1 (doctrine), Deception 1\\
\textbf{Thematic Mechanic:} \textit{Faith's Fragility} \\
When the Hollow Prophet successfully uses Sway to convert a bystander or reinforce a follower's faith, they generate \textbf{1 extra Story Beat} (beyond any generated by 1s). However, if the Sway attempt \textbf{fails} (Miss), they gain a \textit{Despair} Condition: −−1 die to all Spirit and Wits rolls until they successfully persuade someone or gain 1 Boon.\\
\textit{Flavor Hook:} ``Their eyes glow with conviction as they speak−−−but when you catch their reflection in a dark window, you see only exhaustion and doubt. The lie sustains them... but how long until it cracks?''\\
\textit{Rules Tie−In:} Uses Story Beats (\autoref{sec:outcome−sb}) and Conditions (see Chapter~\ref{ch:gm−core}).

⫕section{The Gearsmith}
\textit{Archetype:} Obsessed tinkerer, weaponsmith, or arcane engineer. Sees people as problems to solve.\\
\textit{Stats (Threat Level 2):} Body 2, Wits 3, Spirit 1, Presence 2   |   \textit{Skills:} Craft 2, Mechanics 1, Brawl 1\\
\textbf{Thematic Mechanic:} \textit{Improvised Countermove (Cost: 1 Boon)} \\
When the Gearsmith is targeted by an action (attack, skill use, etc.), the GM may spend 1 Story Beat (or 1 Boon from the GM's pool −− see \autoref{sec:boons−gm}) to invent a reactive device on the spot. The device negates the incoming action and imposes one relevant Condition (e.g., \textit{Entangled}, \textit{Blinded}, \textit{Shaken}) on the attacker. The device is single−use and narratively described (GM improvises based on the scene). This ability can be used at most once per round.\\
\textit{Flavor Hook:} ``They don't dodge your sword−−−they're already calculating how to turn its momentum against you. That clang wasn't a block; it was the sound of a spring−loading.''\\
\textit{Rules Tie−In:} Uses Story Beats (\autoref{sec:outcome−sb}) and Conditions (see Chapter~\ref{ch:gm−core}).

⫕section{The Iron−Bound Rival (PC Mirror)}
\textit{Archetype:} Former ally, rival adventurer, or ideological opponent who mirrors the PCs' path.\\
\textit{Stats:} Match the PC's highest two Attributes at their current rating. Match the PC's two highest Skills at their current rating. Other stats default to 1 or 0 as appropriate.\\
\textbf{Thematic Mechanic:} \textit{Anticipate Move} \\
When a PC declares an action targeting the Rival, the Rival gains \textbf{+1 die} to any defensive or opposing roll (e.g., dodging an attack, resisting Sway, countering a skill use). This bonus does not stack if multiple PCs act against the Rival in the same exchange.\\
\textit{Flavor Hook:} ``They move not like a fighter, but like a shadow of your own choices−−−always one step ahead of where you \textit{think} you're going, but never where you \textit{are}.''\\
\textit{Rules Tie−In:} Uses opposed rolls (see Chapter~\ref{ch:gm−core}) and mirrors PC progression.

⫕section{The Weeping Martyr}
\textit{Archetype:} Someone who has lost everything and fights with suicidal abandon.\\
\textit{Stats (Threat Level 2):} Body 3, Wits 1, Spirit 3, Presence 1   |   \textit{Skills:} Melee 2, Endurance 2, Intimidation (Presence+Brawl)\\
\textbf{Thematic Mechanic:} \textit{Blood Price} \\
When the Martyr takes Harm, they immediately gain +1 die on their next attack or physical action (pain fuels fury). Additionally, if they are reduced to Harm 3 (incapacitated), they may make one final attack as a free action before falling. This attack ignores all Position penalties and deals +1 Harm.\\
\textit{Flavor Hook:} ``They don't even try to block your strike−−−they lean into it, grinning through bloody teeth. Every wound is another coin in a purse they'll never spend.''\\
\textit{Rules Tie−In:} Uses Harm levels (see Chapter~\ref{ch:gm−core}) and the ``roller−coaster'' effect (see \ref{sec:roller−coaster}).

⫕section{The Shadow Broker}
\textit{Archetype:} Information dealer, blackmailer, or spymaster who never fights fair.\\
\textit{Stats (Threat Level 2):} Body 1, Wits 4, Spirit 2, Presence 2   |   \textit{Skills:} Subterfuge 3, Insight 2, Streetwise 2\\
\textbf{Thematic Mechanic:} \textit{I Know a Guy (Cost: 1 Story Beat)} \\
Once per scene, the Shadow Broker may spend 1 Story Beat to declare that they have prepared for the current situation. They produce a minor but useful asset (e.g., a hidden exit, a bribed guard, a forged document) that grants the party (or their own allies) +1 die or improves Position by one step on a single roll. The GM narrates the asset and its limitations.\\
\textit{Flavor Hook:} ``You never see them draw a weapon. You only see the trap they set three scenes ago, the witness they paid off, the escape route you didn't know existed.''\\
\textit{Rules Tie−In:} Uses Story Beats (\autoref{sec:outcome−sb}) and Assets (\autoref{sec:followers−assets−gm}). The Broker rarely engages in direct combat; instead, they create advantages.

\subsection{Using These Templates in Play}
\begin{itemize}
    ℹtem \textbf{Generic foes}: Drop in instantly. Adjust Threat Level via Table~\ref{tab:threat−scaling}. No prep beyond naming them.
    ℹtem \textbf{Archetypal foes}: Pick one matching your scene's theme. Note their one mechanic. Use the \textit{Flavor Hook} to guide narration.
    ℹtem \textbf{Scaling}: If PCs are Tier III, treat Threat Level 2 as baseline. For major encounters, use Threat Level 3 or 4 archetypes.
    ℹtem \textbf{No extra bookkeeping}: Mechanics reference existing rules−−−just apply them as written. The mechanic is the \textit{only} thing to remember beyond core stats.
\end{itemize}

\noindent
\rule{łinewidth}{0.5pt}
\small
\textit{Design Philosophy: Generic foes exist to threaten through numbers and positioning (using core rules only). Archetypal foes add narrative depth via one self−contained mechanic that reflects their core theme−−−never requiring tracking beyond the encounter. All mechanics are verifiable in the core rules (see citations in the text).}

\clearpage
\chapter{Patron Quick Reference}
łabel{chap:patron−quick−ref}

\begin{center}
\begin{longtable}{p{2.2cm}p{2.5cm}p{3.5cm}p{1.6cm}p{3.2cm}}
∩tion{Patron Reference: Themes, Rites, Rivals and Lore}
łabel{tab:patron−quick−ref}\\
⊤rule
\textbf{Patron} & \textbf{Theme} & \textbf{Rites (summary)} & \textbf{Rivals} & \textbf{Misc / Cult / Tool} \\
∣rule
\endfirsthead
μlticolumn{5}{c}%
{{\tablename\ θble{} −− continued}} \\
⊤rule
\textbf{Patron} & \textbf{Theme} & \textbf{Rites (summary)} & \textbf{Rivals} & \textbf{Misc / Cult / Tool} \\
∣rule
\endhead
⊥tomrule
μlticolumn{5}{r}{\textit{Continued on next page}} \\
\endfoot
⊥tomrule
\endlastfoot

Aveh & Freedom & Erasure & Unlatched Step, Forgotten Road, Unremembered Name & — & Ykrul spirit; Cult of the Unbannered; Whistle of the Unbound Road \\
ℎline
Carrion−King & Decay, Renewal & Transformation & Fertile Death, Borrowed Form, Eternal Cycle & — & Compost−Circle; Bone−Needle \\
ℎline
Clockwork Monad & Refinement & Progressive Creation & Anima, self‑optimising devices, transfer inefficiency & — & Iteration Chorus; Brass Stamp of Iteration \\
ℎline
Confessor Beneath the Bell & Burden & Exchange & Absorb sins, transfer conditions, sin‑eater pacts & — & Silent Bell; Bell of Unspoken Burdens \\
ℎline
Gallow’s Bell & Justice & Judgment & Establish accords, weigh hearts, compel truth & — & Unstruck Bell; Seal of Unspoken Weights \\
ℎline
Grimmir & Primal Mystery & Seasonal Sacrifice & Blood offering, speak with roots, consume heart, curse with terrain & — & Root−Moot; Knotted Staff \\
ℎline
Ikasha & Latent Potential & Shadow & Unlit Candle, Crossroads Raven, Umbral Reservoir & — & Unlit Candle; Raven’s Quill \\
ℎline
Inaea & Webs, Fate & Patient Predation & Hearth‑Thread Knot, Snaring Filament, Strand of Inevitability & — & Woven Circle; True Loom \\
ℎline
Inquisitor Prime & Purity & Domination & Resist influence, dispel magic, mark prey, annihilate foes & Oath of Mercy and Grace & Iron Chorus; Censer of the Unforgiving Flame \\
ℎline
Isoka & Transformation & Renewal & Shedding, venomous healing or curse, complete metamorphosis & — & Coiled Heart; Shedding Knife \\
ℎline
Khemesh & Depths, Inexorability, Eldritch Terror & Crushing whispers, bind with force, alien truths, fold reality & Kuva & Drowned Chorus; Vessel of the Unspoken Depth \\
ℎline
Kuva & Sky, Witness, Storm, Threshold of the Hollow & Witness oaths, call lightning, pray for weather shifts & Khemesh, Pale Shepherd & Unclouded Eye; Wind‑Knot Whistle \\
ℎline
Livaea & Seduction & Social Binding & Charming whispers, binding vows, glamour, social sovereignty & Malachai & Velvet Circle; Gossip Quill \\
ℎline
Lunara & Moon, Mystery & Hidden Ways & Mirror scrying, secret thoughts, reveal hidden paths, impossible ways & — & Silver Spiral; Silver Needle \\
ℎline
Mab & Glamour & Bargain & Fae bargains, courtly guises, fae courts, lasting masquerades & — & Gilded Jest; Mask of Shifting Faces \\
ℎline
Maelstraeus & Commerce & Exchange & Trading grounds, appraise value, weighted contracts, cursed market‑zones & — & Gilded Ledger; Debtor’s Quill \\
ℎline
Malachai & Curses & Corruption & Cursed gifts, advantages with corruption timers, supernatural afflictions & Oath of Flame & Light, Livaea & Leashed Hunger; Splintered Needle \\
ℎline
Mor’iraath & Absolute Annihilation & Consume objects, see breaking points, essence‑burning fire, negate existence & — & Final Breath; Cinder‑Stamp \\
ℎline
Morag the Hag & Faerie Bargains & Hidden Costs & Escape clauses, fae knots, hidden‑value trades, crossroads court & — & Tarnished Cauldron; Bargain Needle \\
ℎline
Mykkiel & Law & Zeal & Consecrate areas, divine commandments, covenant penalties, divine judgment & Seal‑Breaker & Gilded Scale; Seal of Unbroken Testimony \\
ℎline
Nidhoggr & Dreaming Antiquity & Learn history from stone, read accumulated memories, ghostly echoes & — & Sleeping Stone; Eye of the World‑Worm \\
ℎline
Nimorith & Absence & Threshold & Create gaps, hold memories as collateral, erase clauses, safe‑conduct, reverse directions & — & Absent Host; Duskmark Lantern \\
ℎline
Ninth (The) & Infinite Information & Unknowable Truths & Glimpse beyond ordinary frames, hypercognitive lattice, loci of overfull knowing & — & Unfinished Index; Paradox Quill \\
ℎline
Oath of Flame & Light & Dawn & Vows & Kindle vow, lay on hands, smite with dawnfire, reveal illusions, covenant zones & Malachai & Unquenched Flame; Prayer Beads of the Unbroken Vow \\
ℎline
Oath of Mercy and Grace & Healing & Forgiveness & Share fatigue healing, remove harm, pacifist vows, sanctuary, unburden debtors & Inquisitor Prime, Malachai & Grey Spiral; Cord of the Unbroken Spiral \\
ℎline
Oya & Dust, Change, Liberation & Extinguish fires, immune to sandstorms, dust storms that reveal truth & Witness, Inquisitor Prime & Red Dust; Dust Pouch \\
ℎline
Pale Shepherd & Passage & Melancholic Wisdom & Ease crossings, find hidden paths, open shadow‑passages, soothe grief & Kuva & Grey Flock; Wool‑Winder \\
ℎline
Palinode & Performance & Rapture & Hymn against dread, perfect performance, captivate audiences, consecrated stages & — & Gasping Chorus; Conductor’s Baton of Unfinished Phrases \\
ℎline
Raéyn & Tides & Change & Stabilise footing, learn sea routes, bias currents, shape storms, curse wrongdoers & Khemesh & Salt‑Washed Circle; Wave‑Seer’s Shell \\
ℎline
Sacred Geometry & Harmony & Immutable Law & Align to harmonics, impose perfect form, harmonic fields, perceive governing patterns & Savage Heart & Unbroken Compass; Protractor of Unseen Measures \\
ℎline
Savage Heart (Red Maw) & Apex Predation & Mark prey, track by scent, assume predator forms, form hunting packs, consume fallen & Grimmir, Sealed Gate, Sacred Geometry & Painted Maw; Hunter’s Mask \\
ℎline
Seal‑Breaker & Forgotten Names, Broken Seals & Present forgeries, strike without traces, permanently void formal contracts & Mykkiel, Witness & Unbound Name; Unpicking Needle \\
ℎline
Sealed Gate & Thresholds & Containment & Seal thresholds, push back intruders, protective circles, banish entities, consecrate barriers & Savage Heart & Unbroken Seal; Seal of the Unbroken Gate \\
ℎline
Silent Choir & Restraint & Protective Omission & Hush spaces, protect secrets, bind unspoken accords, take emotional weight & Witness, Ikasha & Sealed Lip; Candle of Unspoken Trust \\
ℎline
Thrysos & Ecstasy, Excess & Salacious Delight & Loosen inhibitions, ecstatic frenzies, break social masks, summon revellers, endless feast & — & Horned Circle; Thyrsus of the Unbound \\
ℎline
Traveler & Ways & Journeys & Find safe routes, ignore difficult terrain, mark safe paths, relocate targets, bind safe passage & — & Unbroken Way; Way Chisel \\
ℎline
Varnek Karn & Death’s Negotiator & Speak with dead, consult ancestors, mediate pacts, summon Unquiet Host, seal testaments & — & Unforgotten Ledger; Seal of the Silent Witness \\
ℎline
Vorthak & Consumption & Metamorphic Hunger & Absorb traits by touch, permanently consume an item’s nature, adopt aspects of others & — & Ravenous Circle; Maw‑Stone \\
ℎline
Witness & Truth & Revelation & Follow hidden traces, force honest answers, weaken memory and deception, bind agreements, transfer psychological conditions & Seal‑Breaker, Silent Choir & Unmasking (apocalyptic); Quill of Unwritten Truths \\
ℎline
Xhak’Thul & Storm, Shock & Awe & Overawe with presence, move with storm’s speed, call supernatural storms, feed on fear, become incarnation & Traveler, Varnek Karn & First Drum; Storm‑Drum \\
ℎline
\end{longtable}
\end{center}
